\section{CPU}

Centrální procesorová jednotka (CPU) systému NES je implementována prostřednictvím čipu Ricoh R2A03
. Jedná se o~8bitový mikroprocesor odvozený z~architektury MOS Technology 6502, od níž se odlišuje
primárně absencí podpory binárně zakódované dekadické soustavy (BCD režimu) z~důvodu obcházení
patentových práv. Čip R2A03 se vyznačuje tím, že na jednom křemíkovém plátu integruje jak samotné
jádro procesoru, tak i~obvody zvukového generátoru (APU, viz podkapitola \ref{sec:pp-apu}), čímž de
facto představuje ranou formu systému na čipu (SoC).

Pracovní frekvence procesoru je exaktně odvozena dělením hlavního hodinového signálu (master clock).
V~systému NTSC činí hlavní hodinový kmitočet přibližně 21,477 MHz, přičemž procesor pracuje
s~dělitelem 12 na frekvenci 1,789 MHz.

\subsection{Architektura registrů}

Registrová sada procesoru je celkem minimalistická a odpovídá původnímu návrhu architektury 6502.
Procesor nedisponuje univerzálními registry obecného použití, nýbrž využívá sadu úzce
specializovaných registrů:

\begin{itemize}
    \item \textbf{A (Accumulator)} — 8bitový pracovní registr.
    \item \textbf{X, Y (Index Registers)} — dvojice 8bitových indexových registrů.
    \item \textbf{PC (Program Counter)} — 16bitový čítač instrukcí uchovávající paměťovou adresu
        aktuálně zpracovávané instrukce.
    \item \textbf{S (Stack Pointer)} — 8bitový ukazatel vrcholu zásobníku. Hardwarový zásobník je
        fixně alokován na nulté stránce paměti (přesněji v~rozsahu \texttt{\$0100–\$01FF}) a~podle
        konvence roste směrem k~nižším adresám od počáteční hodnoty \texttt{\$01FF}.
    \item \textbf{P (Processor Status)} — 8bitový stavový registr, který sdružuje šest příznaků
        (flags): Carry, Zero, Interrupt Disable, Decimal (v architektuře NES hardwarově
        deaktivován), Overflow a~Negative.
\end{itemize}

\subsection{Paměťová organizace a adresování}

Procesor adresuje lineární 16bitový prostor o~celkové kapacitě 64 KiB. Tento adresní rozsah je
sdílen mezi interní RAM pamětí, paměťově mapovanými periferiemi a~pamětí, určené pro účely mapovače.
Architektura nevyužívá dedikované I/O instrukce; komunikace s~periferiemi probíhá standardními
instrukcemi pro čtení a~zápis do paměti.

\begin{table}[htbp]
    \centering
    \begin{tabular}{lll}
    \toprule
    Rozsah adres & Velikost (byte) & Alokované zařízení \\
    \midrule
    \texttt{\$0000–\$07FF} & 2048 & Interní pracovní RAM \\
    \texttt{\$0800–\$1FFF} & 6144 & Hardwarové zrcadlení interní RAM \\
    \texttt{\$2000–\$2007} & 8 & Mapováné registry grafického koprocesoru (PPU) \\
    \texttt{\$2008–\$3FFF} & 8184 & Zrcadlení PPU registrů (opakující se bloky 8 bytů) \\
    \texttt{\$4000–\$4017} & 24 & Registry APU a~vstupně-výstupní rozhraní (I/O) \\
    \texttt{\$4020–\$FFFF} & 49152 & Prostor cartridge (ROM, RAM, mapovací registry) \\
    \bottomrule
    \end{tabular}
    \caption{Zjednodušená mapa paměťového prostoru CPU}
\end{table}

Interní RAM o~kapacitě 2 KiB je logicky segmentována. Nultá stránka (Zero Page,
\texttt{\$0000–\$00FF}) umožňuje optimalizované adresování pomocí kratších a~rychlejších
dvoubajtových instrukcí. Následuje hardwarově vyhrazená stránka zásobníku (\texttt{\$0100–\$01FF})
a~zbývající kapacita je ponechána pro aplikační data. V~nejvyšších paměťových adresách
(\texttt{\$FFFA–\$FFFF}) jsou pevně definovány vektory přerušení, jež systém načítá z~připojené
cartridge.

\subsection{Přerušení}

Architektura poskytuje následující přerušení:

\begin{itemize}
    \item \textbf{NMI (Non-Maskable Interrupt)} — softwarově nemaskovatelné přerušení, generované
        primárně grafickým procesorem (PPU) při zahájení vertikálního zatemnění (VBlank). Obslužný
        vektor tohoto přerušení se načítá z~adres \texttt{\$FFFA–\$FFFB}.
    \item \textbf{IRQ (Interrupt Request)} — maskovatelné přerušení. Softwarově jej lze potlačit
        nastavením příznaku \textit{Interrupt Disable} ve stavovém registru (P). Obvykle slouží pro
        externí časovače integrované v~mapperu cartridge, nebo pro DMC v APU. Obslužný vektor sídlí
        na adresách \texttt{\$FFFE–\$FFFF}.
    \item \textbf{RESET} — hardwarová inicializace spuštěná při startu konzole nebo po stisku
        tlačítka Reset. Z~hlediska vnitřní logiky CPU funguje jako specifické přerušení s~tím
        rozdílem, že potlačuje instrukce zápisu návratové adresy na zásobník a~automaticky nastavuje
        příznak \textit{Interrupt Disable}. Vektor inicializace se nachází na
        \texttt{\$FFFC–\$FFFD}.
\end{itemize}

\subsection{Adresovací režimy}

Instrukční sada procesoru 6502, z~níž čip R2A03 vychází, specifikuje celkem třináct adresovacích
režimů. Tyto režimy definují způsob, jakým operační kód přistupuje k~operandu, což má přímý dopad na
délku instrukce (v bytech) a~počet strojových taktů nezbytných k~jejímu zpracování. Adresovací
režimy jsou podrobně popsány v datasheetu \cite{6502-datasheet}.

\subsection{Neoficiální instrukce}

Oficiální specifikace architektury 6502 definuje 151 validních instrukčních kódů. Ze zbývajících 105
kombinací v~rámci 8bitového prostoru operačních kódů je většina na fyzickém čipu R2A03 reálně
spustitelná. V~odborné literatuře a~scéně emulátorů se tyto instrukce označují jako
\textit{neoficiální} či \textit{nezdokumentované instrukce} (unofficial/illegal opcodes).

Existence těchto kódů není zamýšlenou funkcí, nýbrž artefaktem způsobu, jakým čip dekóduje
instrukce. Řadič procesoru nevyužívá mikrokód v~tradičním slova smyslu, ale kombinační logiku
řízenou programovatelným logickým polem (PLA — \textit{Programmable Logic Array}) obsahujícím
130\texttimes21 (21 bitů) řádků. Neoficiální operační kódy nečekaně aktivují více řádků PLA
současně. Tím dochází ke spuštění nesouvisejících větví dekodéru a~čip vykoná operaci slučující
chování několika standardních instrukcí.\cite{pagetable-illegal-opcodes} Popis a charakteristika
jednotlivých neoficiálních instrukcí je popsáná v \cite{nmos6510-unintended-opcodes}.
